\chapter{縦断データと線形混合モデル}\label{ch:08}

\begin{goalbox}
\begin{itemize}
\item 反復測定データの平均ベクトルの推定量の分布を導き，時点間の変化量の標準誤差を計算して検定できる．
\item 共分散行列から時点間の相関係数を計算できる．
\item 独立・複合対称・AR(1)・無構造の共分散構造のパラメータ数を数え，平均のパラメータと合わせて AIC で選択できる．
\item 線形混合効果モデル $\bm Y=\bm X\bm\beta+\bm Z\bm\gamma+\bm\varepsilon$ の周辺分散 $\bm V=\bm Z\bm G\bm Z\transp+\bm R$ と一般化最小二乗推定量を導出できる．
\item 変量切片が複合対称構造を生むことを示し，与えられた $\bm V$ から $\bm Z,\bm G$ を構成できる．被験者内相関を無視したときの誤りを説明できる．
\end{itemize}
\end{goalbox}

\begin{exambox}
\begin{itemize}
\item \kakomon{2023}{4}：3時点の反復測定（3変量正規）．$\hat{\bm\mu}$ の分布，2つの共分散構造（行列の形は問題文で与えられる）の AIC，相関行列，変化量の平均と標準誤差，$Z$ 検定．
\item \kakomon{2022}{3}：線形混合モデル．$\bm V=\bm Z\bm G\bm Z\transp+\bm R$，変量効果なしの場合の $\hat{\bm\beta}=(\bm X\transp\bm X)^{-1}\bm X\transp\bm Y$ と不偏性，数値計算，
      ブロック対角の複合対称構造を生む $\bm Z,\bm G$，被験者内相関を考慮する意義．
\item \kakomon{2024}{2}（ベースライン調整の3推定量の比較）と \kakomon{2022}{5}（平均への回帰）は，必要な分布やモデルが問題文で与えられる誘導形式である（\cref{app:F}）．
\end{itemize}
\end{exambox}

\flowline{同一個体の\\反復測定}{時点別平均\\変化量}{多変量正規\\混合モデル}{個体間独立\\共分散構造}{GLS・AIC\\分散の計算}{個体内相関}

\section{問題設定}\label{sec:08-setting}
同一の患者から複数時点で測定した値は，患者固有の水準（個体差）を共有するため正の相関をもつ．
この相関を共分散行列として扱い，(a) 時点別の平均と変化量を推測する，(b) 個体差を変量効果として線形混合モデルで表す．
相関を無視すると，推定量の標準誤差を誤る（被験者レベルの効果では多くの場合過小評価する）．

\section{反復測定データ}\label{sec:08-repeated}
\subsection{平均ベクトルの推定と変化量}\label{subsec:08-mean}
\Hissu\Hinshutsu $N$ 人の $T$ 時点の測定値 $\bm Y_i=(Y_{i1},\dots,Y_{iT})\transp\iid\Normal_T(\bm\mu,\bm\Sigma)$，$\bm\Sigma=(\sigma_{st})$ とする（欠測なし）．
\begin{meidai}[平均ベクトルと変化量（\kakomon{2023}{4}）]\label[meidai]{prop:08-mean}
$\hat{\bm\mu}=\frac1N\sum_{i=1}^N\bm Y_i\sim\Normal_T(\bm\mu,\bm\Sigma/N)$．
時点 $s$ から $t$ への変化量の推定量 $\hat\mu_t-\hat\mu_s=\bm c\transp\hat{\bm\mu}$（$\bm c$ は $t$ 成分 $1$，$s$ 成分 $-1$，他は0）は
\begin{equation}
\hat\mu_t-\hat\mu_s\sim\Normal\Bigl(\mu_t-\mu_s,\ \frac{\sigma_{tt}+\sigma_{ss}-2\sigma_{st}}{N}\Bigr).\label{eq:08-change}
\end{equation}
\end{meidai}
\begin{proof}
独立な多変量正規ベクトルの和は多変量正規で，$\E[\hat{\bm\mu}]=\bm\mu$，$\V[\hat{\bm\mu}]=N\bm\Sigma/N^2=\bm\Sigma/N$．
線形変換 $\bm c\transp\hat{\bm\mu}$ の分散は $\bm c\transp\bm\Sigma\bm c/N=(\sigma_{tt}+\sigma_{ss}-2\sigma_{st})/N$．
\end{proof}
$H_0:\mu_t=\mu_s$ は $Z=(\hat\mu_t-\hat\mu_s)/\SE$ で検定する．$\bm\Sigma$ を推定値で置き換えるので $Z$ は大標本で近似的に $\Normal(0,1)$ に従う
（各人の変化量 $Y_{it}-Y_{is}$ に対応のある $t$ 検定（\cref{sec:03-t}）を行うのと同じ推定量である）．
正の相関 $\sigma_{st}>0$ があると，変化量の分散は独立の場合の $(\sigma_{tt}+\sigma_{ss})/N$ より小さくなる．
時点 $s,t$ の相関係数は $\sigma_{st}/\sqrt{\sigma_{ss}\sigma_{tt}}$ である．

\subsection{共分散構造と AIC}\label{subsec:08-cov}
\begin{table}[htbp]
\centering\small
\caption{主な共分散構造（$T$ 時点）}\label{tab:08-cov}
\begin{tabular}{@{}lL{62mm}c@{}}
\toprule
構造 & $(s,t)$ 成分 & パラメータ数\\
\midrule
独立 & $\sigma^2\ind(s=t)$ & 1\\
複合対称（CS，交換可能） & $\sigma_b^2+\sigma^2\ind(s=t)$ & 2\\
AR(1) & $\sigma^2\phi^{|s-t|}$ & 2\\
無構造（UN） & $\sigma_{st}$（すべて自由，$\sigma_{st}=\sigma_{ts}$） & $T(T+1)/2$\\
\bottomrule
\end{tabular}
\end{table}
\Hinshutsu 構造の選択には AIC $=-2\ell+2k$ を使い（\cref{def:06-aic}），$k$ は\textbf{平均のパラメータ数と共分散のパラメータ数の和}である．
問題文で行列の形が与えられた構造では，行列に現れる異なる未知パラメータの個数を数えればよい．
\kakomon{2023}{4}（$T=3$，平均は各時点で自由で3個）では，問題文で与えられた構造（共分散パラメータ3個）が計6個，無構造が $3+6=9$ 個であり，
$\mathrm{AIC}=412.84+12=424.84$ 対 $409.74+18=427.74$ で前者が選ばれる．ただし差は2.9と小さく，決定的とは言えない．
\begin{supbox}[補足：REML]
分散成分の最尤推定量は平均の推定による自由度の損失を考慮しないため下方に偏る（$\hat\sigma^2=\frac1n\sum(x_i-\bar x)^2$ が $\sigma^2$ を過小推定するのと同じ現象）．
混合モデルの実務ではこれを補正した制限付き最尤法（REML）がよく用いられる．REML の尤度は固定効果の構造に依存するので，固定効果が異なるモデルの比較を REML 尤度の AIC で行ってはいけない．
\end{supbox}

\section{線形混合効果モデル}\label{sec:08-lmm}
\subsection{モデル}\label{subsec:08-lmmdef}
\begin{teigi}[線形混合効果モデル]\label{def:08-lmm}
\begin{equation}
\bm Y=\bm X\bm\beta+\bm Z\bm\gamma+\bm\varepsilon,\qquad \bm\gamma\sim\Normal(\bm0,\bm G),\quad \bm\varepsilon\sim\Normal(\bm0,\bm R),\quad \bm\gamma\indep\bm\varepsilon.\label{eq:08-lmm}
\end{equation}
$\bm\beta$ は\textbf{固定効果}（集団平均の構造），$\bm\gamma$ は\textbf{変量効果}（個体ごとの確率的なずれ），$\bm X,\bm Z$ は既知の計画行列．
\end{teigi}
\begin{meidai}[周辺分布（\kakomon{2022}{3}〔1〕）]\label[meidai]{prop:08-marginal}
$\E[\bm Y]=\bm X\bm\beta$，$\V[\bm Y]=\bm V=\bm Z\bm G\bm Z\transp+\bm R$．$\bm Y$ は多変量正規分布 $\Normal(\bm X\bm\beta,\bm V)$ に従う．
\end{meidai}
\begin{proof}
$\bm X\bm\beta$ は定数，$\bm\gamma$ と $\bm\varepsilon$ は独立なので
$\V[\bm Y]=\V[\bm Z\bm\gamma]+\V[\bm\varepsilon]=\bm Z\bm G\bm Z\transp+\bm R$．独立な正規ベクトルの線形結合なので正規．
\end{proof}

\subsection{一般化最小二乗推定}\label{subsec:08-gls}
\begin{teiri}[GLS 推定量]\label[teiri]{thm:08-gls}
$\bm V$ が既知で $\bm X$ が列フルランクなら，$\bm\beta$ の最尤推定量（＝一般化最小二乗推定量）は
\begin{equation}
\hat{\bm\beta}=(\bm X\transp\bm V^{-1}\bm X)^{-1}\bm X\transp\bm V^{-1}\bm Y,\qquad
\E[\hat{\bm\beta}]=\bm\beta,\qquad \V[\hat{\bm\beta}]=(\bm X\transp\bm V^{-1}\bm X)^{-1}.\label{eq:08-gls}
\end{equation}
\end{teiri}
\begin{proof}
対数尤度の $\bm\beta$ を含む部分は $-\frac12(\bm Y-\bm X\bm\beta)\transp\bm V^{-1}(\bm Y-\bm X\bm\beta)$ で，これを最大化すると正規方程式
$\bm X\transp\bm V^{-1}\bm X\bm\beta=\bm X\transp\bm V^{-1}\bm Y$．$\bm A=(\bm X\transp\bm V^{-1}\bm X)^{-1}\bm X\transp\bm V^{-1}$ とおくと
$\E[\hat{\bm\beta}]=\bm A\bm X\bm\beta=\bm\beta$，$\V[\hat{\bm\beta}]=\bm A\bm V\bm A\transp=(\bm X\transp\bm V^{-1}\bm X)^{-1}$．
\end{proof}
実際は $\bm V$ を推定値 $\hat{\bm V}$ で置き換える．
変量効果がなく $\bm R=\sigma^2\bm I$ なら $\hat{\bm\beta}=(\bm X\transp\bm X)^{-1}\bm X\transp\bm Y$（通常の最小二乗；\kakomon{2022}{3}〔2〕）．
このとき $\E[\hat{\bm\beta}]=\bm\beta$ は $\V[\bm Y]$ の形によらず成り立つ（不偏性には $\E[\bm Y]=\bm X\bm\beta$ だけを使う）が，$\V[\hat{\bm\beta}]=\sigma^2(\bm X\transp\bm X)^{-1}$ は相関があると成り立たない．

\subsection{変量切片と複合対称}\label{subsec:08-ri}
\Hissu 被験者 $i$ の $j$ 回目の測定を $Y_{ij}=\bm x_{ij}\transp\bm\beta+b_i+\varepsilon_{ij}$，$b_i\sim\Normal(0,\sigma_b^2)$，$\varepsilon_{ij}\sim\Normal(0,\sigma^2)$（すべて独立）とする（変量切片モデル）．
被験者ごとのブロックで $\bm Z_i=\bm 1_{m}$（$m$ は測定回数），$\bm G_i=\sigma_b^2$，$\bm R_i=\sigma^2\bm I_m$ だから
\[
\bm V_i=\sigma_b^2\bm 1\bm 1\transp+\sigma^2\bm I=
\begin{pmatrix}\sigma^2+\sigma_b^2&\sigma_b^2&\cdots\\ \sigma_b^2&\sigma^2+\sigma_b^2&\cdots\\ \vdots&&\ddots\end{pmatrix},
\qquad \Corr(Y_{ij},Y_{ik})=\frac{\sigma_b^2}{\sigma_b^2+\sigma^2}\ (j\ne k).
\]
これが複合対称構造であり，相関（級内相関係数，ICC）は時点差によらない．
被験者全体では $\bm Z$ は被験者の指示行列（各行は，その測定がどの被験者のものかを示す0/1），$\bm G=\sigma_b^2\bm I$ で，$\bm V$ は $\bm V_i$ を並べたブロック対角行列になる．
\kakomon{2022}{3}〔3〕のように被験者ごとに異なる変量切片の分散を許すなら，$\bm G$ を対角成分が被験者ごとに異なる対角行列にすればよい．
変量傾きモデル $Y_{ij}=\bm x_{ij}\transp\bm\beta+b_{0i}+b_{1i}t_j+\varepsilon_{ij}$ では $\bm Z_i=(\bm1,\bm t)$，$\bm G_i$ は $2\times2$ で，分散が時間とともに変化する構造になる．

\paragraph{相関を無視したときの誤り}
各被験者 $m$ 回測定の変量切片モデルで，被験者レベルの共変量（治療群）の効果を群平均の差で推定するとき，
各群 $n$ 人の群平均の分散は $\frac{\sigma_b^2+\sigma^2/m}{n}=\frac{(\sigma_b^2+\sigma^2)\{1+(m-1)\mathrm{ICC}\}}{nm}$ である．
全観測を独立とみなすと分散を $(\sigma_b^2+\sigma^2)/(nm)$ とするので，真の分散の $1/\{1+(m-1)\mathrm{ICC}\}$ に過小評価する（$1+(m-1)\mathrm{ICC}$ をデザイン効果という）．
変量効果を入れたモデル（\kakomon{2022}{3}〔3〕）は，この被験者内相関を正しく考慮した推測を可能にする．

\subsection{固定効果と変量効果}\label{subsec:08-fixrand}
\begin{itemize}
\item \textbf{固定効果}：関心のある特定の水準の効果（治療，時点，性別）．水準の効果そのものを推定・比較する．
\item \textbf{変量効果}：母集団から抽出されたとみなす水準（被験者，施設）の効果．個々の値ではなく分散 $\sigma_b^2$ に関心がある．
\end{itemize}

\section{仮定}\label{sec:08-assump}
\begin{itemize}
\item 個体間の独立性（個体内の相関は共分散構造でモデル化）．
\item 多変量正規性，平均構造と共分散構造の正しい特定．
\item 欠測がない（\cref{prop:08-mean}）．欠測がある場合，尤度に基づく解析は欠測が MAR なら妥当，完全ケース解析は MCAR を要する（\cref{ch:16}）．
\end{itemize}

\section{結果の解釈}\label{sec:08-interp}
\begin{itemize}
\item 変化量の検定が有意でも，対照群のない単群の前後比較では時間経過や測定の変動による変化と処置の効果は区別できない．
\item 共分散構造の選択は推定値そのものより標準誤差に影響する．AIC の差が小さいときは複数の構造で結論が変わらないかを確認する．
\item ICC が大きいほど，同じ被験者の追加測定から得られる情報は少ない（デザイン効果が大きい）．
\end{itemize}

\section{手法選択}\label{sec:08-choice}
\begin{itemize}
\item 欠測のない多時点の連続データで時点間の変化 → 平均ベクトルの推定と変化量の $Z$ 検定（\cref{eq:08-change}）．
\item 共分散構造の選択 → 平均と共分散のパラメータ数を数えて AIC．
\item 個体差を含めて平均構造を推定 → 線形混合効果モデル（変量切片なら複合対称）と GLS 推定．
\item 被験者レベルの比較で全観測を独立とみなす解析 → 誤り（デザイン効果の分だけ分散を過小評価）．
\end{itemize}

\section{よくある誤り}\label{sec:08-err}
\begin{warnbox}
\begin{itemize}
\item 変化量の分散で共分散の項 $-2\sigma_{st}$ を落とす．
\item 反復測定の全観測を独立とみなして $t$ 検定・OLS の標準誤差を使う．
\item 無構造のパラメータ数を $T^2$ とする（正しくは $T(T+1)/2$）．AIC で平均のパラメータを数え忘れる．
\item $\hat{\bm\mu}$ の共分散行列を $\bm\Sigma$ とする（正しくは $\bm\Sigma/N$）．
\item 変量切片モデルの相関を時点差によって変わるとする（複合対称なので一定）．
\end{itemize}
\end{warnbox}

\section{数理統計との接続}\label{sec:08-link}
\begin{linkbox}
\begin{itemize}
\item 多変量正規の線形変換（『現代数理統計学の基礎』命題4.25）から $\hat{\bm\mu}$ と GLS 推定量の分布が得られる．
\item 条件付き分散・共分散の公式（同書 命題4.11，第4章演習）で $\V[\bm Y]=\E[\V[\bm Y\mid\bm\gamma]]+\V[\E[\bm Y\mid\bm\gamma]]=\bm R+\bm Z\bm G\bm Z\transp$ とも導ける．
\item 分散の MLE の除数 $n$ による偏り（同書 式(5.5)，\kakomonSM{2017}{1}）が REML の動機である．
\item \kakomonSM{2022}{5} は対応のある比較を二元配置分散分析として扱い，被験者を固定効果とみるか変量効果とみるかの考え方に通じる．
\end{itemize}
\end{linkbox}

\section{例題}\label{sec:08-ex}
\begin{reidai}[反復測定]\label{reidai:08-1}
30人の3時点（0，4，8週）の測定値の平均ベクトルの推定値が $(10.2,9.1,8.3)$，
共分散行列の推定値が $\begin{pmatrix}4.0&2.4&2.0\\2.4&4.4&2.6\\2.0&2.6&4.8\end{pmatrix}$ である．
\begin{enumerate}[label=(\arabic*)]
\item 0週から8週の変化量の推定値と標準誤差を求め，$H_0:\mu_8=\mu_0$ を有意水準5\%で検定せよ（$z_{0.025}=1.96$）．
\item 0週と8週の相関係数を求めよ（$\sqrt{19.2}=4.382$）．
\item 4時点の研究で，独立，CS，AR(1)，無構造の共分散パラメータ数を答えよ．
\item 3時点で，無構造（最大対数尤度 $-150.2$）と AR(1)（$-153.9$）の AIC を比較せよ．平均は各時点で自由とする．
\end{enumerate}
\end{reidai}

\begin{reidai}[変量切片モデル]\label{reidai:08-2}
$Y_{ij}=\mu+\delta z_i+b_i+\varepsilon_{ij}$（$j=1,\dots,m$，$b_i\sim\Normal(0,\sigma_b^2)$，$\varepsilon_{ij}\sim\Normal(0,\sigma^2)$，独立）で，各群 $n$ 人とする（$z_i\in\{0,1\}$ は群）．
\begin{enumerate}[label=(\arabic*)]
\item 被験者 $i$ の $m$ 個の測定値の共分散行列を書き，これを $\bm Z_i\bm G_i\bm Z_i\transp+\bm R_i$ の形に表せ．
\item 群平均の差 $\hat\delta$ の分散を求め，全観測を独立とみなした場合の分散との比（デザイン効果）を求めよ．
\item $\sigma_b^2=3$，$\sigma^2=1$，$m=4$ のときデザイン効果はいくらか．
\end{enumerate}
\end{reidai}

\section{解答}\label{sec:08-sol}
\begin{kaitou}{reidai:08-1}
(1) $8.3-10.2=-1.9$．$\widehat\V=(4.0+4.8-2\times2.0)/30=4.8/30=0.16$，$\SE=0.40$．$Z=-1.9/0.40=-4.75$，$|Z|>1.96$ で有意（共分散行列を推定値で置き換えているので，$Z$ は大標本で近似的に $\Normal(0,1)$）．\\
(2) $2.0/\sqrt{4.0\times4.8}=2.0/4.382=0.456$．\\
(3) 独立 1，CS 2，AR(1) 2，無構造 $4\cdot5/2=10$．\\
(4) 平均のパラメータは3個．無構造：$300.4+2\times(3+6)=318.4$．AR(1)：$307.8+2\times(3+2)=317.8$．AR(1) の AIC が小さい（ただし差は0.6でわずか）．
\end{kaitou}

\begin{kaitou}{reidai:08-2}
(1) 対角成分 $\sigma_b^2+\sigma^2$，非対角成分 $\sigma_b^2$ の $m\times m$ 行列（複合対称）．$\bm Z_i=\bm1_m$，$\bm G_i=\sigma_b^2$，$\bm R_i=\sigma^2\bm I_m$ として $\sigma_b^2\bm1\bm1\transp+\sigma^2\bm I_m$．\\
(2) 被験者平均 $\bar Y_{i\cdot}=\mu+\delta z_i+b_i+\bar\varepsilon_{i\cdot}$ の分散は $\sigma_b^2+\sigma^2/m$ で，群平均はその $n$ 人の平均，2群は独立なので
$\V[\hat\delta]=2(\sigma_b^2+\sigma^2/m)/n$．独立とみなすと $2(\sigma_b^2+\sigma^2)/(nm)$．
比は $\dfrac{m\sigma_b^2+\sigma^2}{\sigma_b^2+\sigma^2}=1+(m-1)\mathrm{ICC}$，$\mathrm{ICC}=\sigma_b^2/(\sigma_b^2+\sigma^2)$．\\
(3) $\mathrm{ICC}=3/4$，$1+3\times0.75=3.25$．独立を仮定すると分散を約3分の1に過小評価する．
\end{kaitou}

\begin{summarybox}
\begin{itemize}
\item $\hat{\bm\mu}\sim\Normal(\bm\mu,\bm\Sigma/N)$，変化量の分散 $(\sigma_{tt}+\sigma_{ss}-2\sigma_{st})/N$，相関 $\sigma_{st}/\sqrt{\sigma_{ss}\sigma_{tt}}$．
\item 共分散構造のパラメータ数：独立 1，CS 2，AR(1) 2，UN $T(T+1)/2$．AIC は平均＋共分散のパラメータ数で計算．
\item LMM：$\bm V=\bm Z\bm G\bm Z\transp+\bm R$，$\hat{\bm\beta}=(\bm X\transp\bm V^{-1}\bm X)^{-1}\bm X\transp\bm V^{-1}\bm Y$，$\V[\hat{\bm\beta}]=(\bm X\transp\bm V^{-1}\bm X)^{-1}$．
\item 変量切片 → 複合対称，ICC $=\sigma_b^2/(\sigma_b^2+\sigma^2)$．相関を無視するとデザイン効果 $1+(m-1)\mathrm{ICC}$ だけ分散を過小評価．
\end{itemize}
\end{summarybox}
